\vspace{0.5em}\noindent\textbf{Tổng quan}\par

Workshop này trình bày toàn bộ quy trình triển khai ứng dụng Internship
Application Tracker trên nền tảng AWS.

Nội dung workshop bao gồm chuẩn bị mã nguồn, khởi tạo hạ tầng cloud,
triển khai cơ sở dữ liệu, triển khai ứng dụng, xử lý AI, gửi thông báo
theo sự kiện, giám sát, kiểm thử, xử lý sự cố và dọn dẹp tài nguyên.

Kiến trúc chính thức sử dụng:

\begin{itemize}
\tightlist
\item
  Amazon CloudFront và Amazon S3 để phân phối frontend.
\item
  Application Load Balancer để tiếp nhận lưu lượng public API và thời
  gian thực (realtime).
\item
  Amazon EKS cho các workload backend, chat, worker và AI adapter.
\item
  Amazon RDS PostgreSQL để lưu trữ dữ liệu giao dịch nghiệp vụ.
\item
  Amazon ElastiCache Redis phục vụ kênh pub/sub thời gian thực cho
  Socket.IO.
\item
  Amazon DynamoDB để lưu trữ bền vững dữ liệu chat và khử lặp
  (deduplication) sự kiện cho Lambda.
\item
  Amazon SQS để truyền tải sự kiện bất đồng bộ.
\item
  AWS Lambda và Amazon SES để gửi thông báo.
\item
  Amazon SageMaker để phục vụ suy đoán AI (AI inference).
\item
  GitHub Actions và AWS OIDC cho luồng triển khai tự động liên tục.
\end{itemize}

\vspace{0.5em}\noindent\textbf{Mục tiêu học tập}\par

Sau khi hoàn thành workshop này, người đọc có thể:

\begin{enumerate}
\def\labelenumi{\arabic{enumi}.}
\tightlist
\item
  Chuẩn bị một ứng dụng kiến trúc đa dịch vụ (multi-service) sẵn sàng để
  triển khai lên AWS.
\item
  Triển khai các dịch vụ container lên Amazon EKS.
\item
  Triển khai frontend tĩnh React thông qua Amazon S3 và CloudFront.
\item
  Kết nối các workload trên EKS với RDS, Redis, DynamoDB, SQS và
  SageMaker.
\item
  Thực hiện xử lý bất đồng bộ bằng các worker và mô hình transactional
  outbox.
\item
  Xử lý các sự kiện SQS bằng AWS Lambda.
\item
  Triển khai cơ chế xử lý tính bất biến (idempotent) cho Lambda bằng
  cách sử dụng DynamoDB.
\item
  Cấu hình luồng CI/CD sử dụng GitHub Actions và AWS OIDC.
\item
  Kiểm chứng toàn bộ hệ thống thông qua kiểm thử đầu cuối (end-to-end
  testing).
\item
  Chẩn đoán các lỗi mạng, IAM và triển khai phổ biến.
\end{enumerate}

\vspace{0.5em}\noindent\textbf{Đối tượng hướng
đến}\par

Workshop này phù hợp dành cho:

\begin{itemize}
\tightlist
\item
  Thực tập sinh kỹ thuật Cloud (Cloud engineering interns).
\item
  Sinh viên theo học hướng Backend và DevOps.
\item
  Lập trình viên đang tìm hiểu và thực hành với Amazon EKS.
\item
  Lập trình viên muốn xây dựng kiến trúc hướng sự kiện (event-driven
  architectures) trên AWS.
\item
  Sinh viên đang thực hiện báo cáo thực tập hoặc đồ án tốt nghiệp về hệ
  thống AWS.
\end{itemize}

\vspace{0.5em}\noindent\textbf{Tóm tắt kiến trúc}\par

\begin{figure}
\centering
\pandocbounded{\includegraphics[keepaspectratio,alt={Kiến trúc hạ tầng Internship Application Tracker trên AWS}]{5-Workshop/infra.png}}
\caption{Kiến trúc hạ tầng Internship Application Tracker trên AWS}
\end{figure}

Trình duyệt người dùng kết nối an toàn tới Amazon CloudFront thông qua
giao thức HTTPS.

CloudFront định tuyến:

\begin{itemize}
\tightlist
\item
  các request tải trang tĩnh của frontend đến Amazon S3
\item
  các request \texttt{/\allowbreak{}api/\allowbreak{}*} đến dịch vụ backend thông qua Application
  Load Balancer
\item
  các request \texttt{/\allowbreak{}chat/\allowbreak{}*} đến dịch vụ chat
\item
  các request \texttt{/\allowbreak{}socket.\allowbreak{}io/\allowbreak{}*} đến dịch vụ liên lạc thời gian thực
  Socket.IO
\end{itemize}

Backend, dịch vụ chat, processing worker, dịch vụ AI và outbox
dispatcher vận hành bên trong Amazon EKS.

PostgreSQL lưu trữ dữ liệu giao dịch nghiệp vụ và các tác vụ processing
jobs. DynamoDB lưu trữ dữ liệu tin nhắn chat. Redis cung cấp cơ chế
pub/sub điều phối tin nhắn giữa các pod chat.

Processing worker gọi đến dịch vụ AI (AI service), và dịch vụ AI sẽ trực
tiếp gửi lời gọi đến máy chủ SageMaker endpoint.

Outbox dispatcher tiến hành phát hành các sự kiện nghiệp vụ sang Amazon
SQS. AWS Lambda tiếp nhận các sự kiện này, thực thi logic khử lặp trong
DynamoDB, lưu giữ (archive) nội dung sự kiện vào Amazon S3 và phát tán
email thông báo thông qua Amazon SES.

\vspace{0.5em}\noindent\textbf{Danh sách các
chương}\par

{
\begingroup
\small
\setlength{\tabcolsep}{3pt}
\renewcommand{\arraystretch}{1.15}
\begin{longtable}{@{}
  >{\raggedright\arraybackslash}p{(\linewidth - 6\tabcolsep) * \real{0.2500}}
  >{\raggedright\arraybackslash}p{(\linewidth - 6\tabcolsep) * \real{0.2500}}
  >{\raggedright\arraybackslash}p{(\linewidth - 6\tabcolsep) * \real{0.2500}}
  >{\raggedright\arraybackslash}p{(\linewidth - 6\tabcolsep) * \real{0.2500}}@{}}
\toprule\noalign{}
\begin{minipage}[b]{\linewidth}\raggedright
Chương
\end{minipage} & \begin{minipage}[b]{\linewidth}\raggedright
Tiêu đề
\end{minipage} & \begin{minipage}[b]{\linewidth}\raggedright
Mô tả
\end{minipage} & \begin{minipage}[b]{\linewidth}\raggedright
Trạng thái
\end{minipage} \\
\midrule\noalign{}
\endhead
\bottomrule\noalign{}
\endlastfoot
5.1 & Tổng quan & Mục tiêu, phạm vi và kiến trúc hoàn chỉnh của workshop
& Hoàn thành \\
5.2 & Kiến trúc & Chi tiết kiến trúc AWS và luồng di chuyển của các
request & Hoàn thành \\
5.3 & Điều kiện tiên quyết & Yêu cầu về tài khoản, công cụ kỹ thuật và
cấu trúc máy địa phương & Hoàn thành \\
5.4 & Chuẩn bị mã nguồn & Cấu trúc repository, biến môi trường và quy
tắc kiểm định & Hoàn thành \\
5.5 & Hạ tầng AWS & Mạng kết nối, IAM, EKS, ECR, ALB và CloudFront &
Hoàn thành \\
5.6 & Triển khai cơ sở dữ liệu & Triển khai PostgreSQL, Redis và
DynamoDB & Hoàn thành \\
5.7 & Triển khai backend & Triển khai backend, chat, dispatcher, dịch vụ
AI và worker & Hoàn thành \\
5.8 & Triển khai frontend & Build mã React, tải tệp lên S3 và thiết lập
CDN CloudFront & Hoàn thành \\
5.9 & Giám sát và nhật ký & Nhật ký, điểm kiểm tra sức khỏe, thông số
metrics và cảnh báo & Hoàn thành \\
5.10 & Kiểm thử đầu cuối & Nghiệm thu tổng thể ứng dụng, chat, dịch vụ
AI và Lambda & Hoàn thành \\
5.11 & Bảo mật và chi phí & IAM, mã hóa dữ liệu, OIDC và các giả định về
cước phí & Hoàn thành \\
5.12 & Xử lý sự cố & Phân tích và tháo gỡ các lỗi triển khai cũng như
lỗi vận hành thực tế & Hoàn thành \\
5.13 & Dọn dẹp tài nguyên & Quy trình từng bước hủy bỏ an toàn các tài
nguyên mồi gánh chi phí & Hoàn thành \\
\end{longtable}
\endgroup
}

\vspace{0.5em}\noindent\textbf{Kết quả đầu ra}\par

Tại điểm kết thúc workshop, ứng dụng Internship Application Tracker sẵn
sàng đón nhận kết nối qua mạng CloudFront và đảm bảo trọn vẹn các năng
lực:

\begin{itemize}
\tightlist
\item
  phân phối trang tĩnh cho frontend
\item
  hệ thống API từ backend
\item
  liên lạc chat thời gian thực
\item
  lưu trữ dữ liệu an toàn và lâu dài
\item
  tác nghiệp suy đoán AI bất đồng bộ
\item
  khâu phát tin nhắn theo mô hình transactional outbox
\item
  gửi thông tin cảnh báo thông qua AWS Lambda
\item
  luồng triển khai tự động hóa mượt mà nhờ GitHub Actions
\end{itemize}
